\documentclass[aspectratio=169]{beamer}
\usetheme{Madrid}
\usecolortheme{default}

\usepackage{graphicx}
\usepackage{amsmath, amssymb, bm}
\usepackage{hyperref}
\usepackage{xcolor}
\usepackage{booktabs}
\usepackage[backend=biber, style=numeric, sorting=none]{biblatex}
\addbibresource{../docs/source/refs.bib}

\graphicspath{
  {../data/final/orszag_tang_128/figs/}
  {../data/final/convergence/figs/}
  {../data/final/mesh_preview/figs/}
  {../docs/source/figs/}
}

\title[2D Ideal MHD]{2D Ideal MHD Solver on the Orszag-Tang Vortex}
\subtitle{ME-7751 Final Project}
\author{Austin Benny}
\institute{}
\date{\today}

\begin{document}

\begin{frame}
  \titlepage
\end{frame}

\begin{frame}{Outline}
  \tableofcontents
\end{frame}

\section{Introduction and Problem Motivation}

\begin{frame}{Why MHD}
  \begin{itemize}
    \item Magnetohydrodynamics couples fluid flow and magnetic field through the Lorentz force~\cite{chen2016,freidberg2014}.
    \item Valid when plasma is collisional and magnetic Reynolds number is large:
          solar wind, tokamaks, astrophysical jets, magnetospheres.
    \item New physics vs Euler: Alfven and fast/slow magnetosonic waves,
          the Maxwell stress, and the constraint $\nabla\cdot B = 0$.
  \end{itemize}
  \vspace{0.5em}
  \textbf{Project goal.} Build a 2D ideal MHD finite-volume solver in Python and validate it on the Orszag-Tang vortex.
\end{frame}

\begin{frame}{Why Orszag-Tang}
  \begin{columns}[T]
    \begin{column}{0.55\textwidth}
      \begin{itemize}
        \item Canonical 2D MHD benchmark.
        \item All-sinusoidal ICs on $[0,1]^{2}$, periodic on every side.
        \item No geometry, no density jump, no outflow BCs.
        \item If the result is wrong, it is the solver, not the setup.
        \item Reference: \cite{londrillo2000,athena_orszag_tang}.
      \end{itemize}
    \end{column}
    \begin{column}{0.45\textwidth}
      \includegraphics[width=\linewidth]{rho_final.pdf}
    \end{column}
  \end{columns}
\end{frame}

\section{Mathematical Model}

\begin{frame}{Ideal MHD in conservative form}
  \begin{equation*}
  \frac{\partial U}{\partial t} + \frac{\partial F}{\partial x} + \frac{\partial G}{\partial y} = 0,
  \quad
  U = (\rho,\; \rho\mathbf{v},\; \mathbf{B},\; E)^{T}
  \end{equation*}
  with total energy
  $E = p/(\gamma-1) + \tfrac{1}{2}\rho|\mathbf{v}|^{2} + \tfrac{1}{2}|\mathbf{B}|^{2}$
  and total pressure $p_T = p + \tfrac{1}{2}|\mathbf{B}|^{2}$~\cite{powell1999}.
  \begin{itemize}
    \item Eight conserved variables; all three velocity and field components are carried in 2D.
    \item Seven-wave fan: two fast magnetosonic, two Alfven, two slow magnetosonic, one entropy~\cite{brio1988,miyoshi2005}.
    \item Ideal-gas closure, $\gamma = 5/3$.
  \end{itemize}
\end{frame}

\begin{frame}{The divergence constraint}
  \begin{itemize}
    \item Exact ideal MHD preserves $\nabla\cdot B = 0$ analytically.
    \item A piecewise-constant Godunov update does not: centered finite-differences produce spurious divergence at cell centers~\cite{balsara1999,powell1999}.
    \item Consequences: parallel-to-B Lorentz force, loss of momentum conservation, eventual blow-up.
    \item Remedies: constrained transport~\cite{balsara1999,londrillo2000}, Powell source terms~\cite{powell1999}, or hyperbolic GLM cleaning~\cite{dedner2002}.
    \item \textbf{This project uses GLM}: simpler to drop into an unstaggered FVM code, only adds one scalar $\psi$.
  \end{itemize}
\end{frame}

\section{Numerical Model}

\begin{frame}{Computational mesh}
  \begin{itemize}\scriptsize
    \item Uniform $128\times 128$ Cartesian grid, $\Delta x = \Delta y = 1/128$.
    \item Two ghost layers on every side (periodic fills for Orszag-Tang).
    \item Left: initial velocity; right: initial magnetic field.
  \end{itemize}
  \vspace{0.2em}
  \centering
  \includegraphics[width=0.78\textwidth]{mesh_preview.pdf}
\end{frame}

\begin{frame}{Scheme at a glance}
  \begin{itemize}
    \item 2D structured Cartesian grid, cell-centered state, two ghost layers~\cite{versteeg2007,moukalled2016}.
    \item MUSCL linear reconstruction on primitives with minmod limiter (MC, van Leer optional)~\cite{toro2009}.
    \item HLLD Riemann solver~\cite{miyoshi2005} at every interface; degenerates to HLLC for B = 0.
    \item GLM psi/Bx coupling handled inside the Riemann solver; parabolic damping applied per RK stage~\cite{dedner2002}.
    \item SSP-RK2 time integration~\cite{shu1988}, CFL = 0.3.
  \end{itemize}
\end{frame}

\begin{frame}{HLLD wave pattern}
  \begin{columns}[T]
    \begin{column}{0.55\textwidth}
      \begin{itemize}
        \item Five-wave approximation to the full seven-wave MHD fan~\cite{miyoshi2005}.
        \item Bounding speeds $S_L, S_R$: Davis-Einfeldt.
        \item Entropy speed $S_M$ from jump conditions at the contact wave.
        \item Two starred states and two double-starred states (across the rotational Alfven discontinuities).
        \item Positivity-preserving for physical inputs.
      \end{itemize}
    \end{column}
    \begin{column}{0.45\textwidth}
      \centering
      \resizebox{\linewidth}{!}{%
      \begin{tabular}{|c|c|c|c|c|c|}\hline
        L & $L^{*}$ & $L^{**}$ & $R^{**}$ & $R^{*}$ & R \\\hline
        \multicolumn{2}{|c|}{$S_L$} & \multicolumn{2}{c|}{$S_L^{*}$} & $S_M$ & $S_R$ \\\hline
      \end{tabular}}
      \vspace{0.5em}
      \scriptsize Wave regions across the x-face.
    \end{column}
  \end{columns}
\end{frame}

\section{Validation and Results}

\begin{frame}{Density at t = 0.5, $128^{2}$}
  \begin{columns}[c]
    \begin{column}{0.55\textwidth}
      \centering
      \includegraphics[width=\linewidth]{rho_final.pdf}
    \end{column}
    \begin{column}{0.45\textwidth}
      \begin{itemize}\small
        \item Central bow-tie current sheet.
        \item Four symmetric oblique shocks bounding low-density bays.
        \item Qualitative match with \cite{londrillo2000,athena_orszag_tang} at the same time slice.
      \end{itemize}
    \end{column}
  \end{columns}
\end{frame}

\begin{frame}{Thermal and magnetic pressure, t = 0.5}
  \begin{columns}[c]
    \begin{column}{0.5\textwidth}
      \centering
      \includegraphics[width=\linewidth]{p_final.pdf}
    \end{column}
    \begin{column}{0.5\textwidth}
      \centering
      \includegraphics[width=\linewidth]{pmag_final.pdf}
    \end{column}
  \end{columns}
  \vspace{0.3em}
  \begin{itemize}\scriptsize
    \item Thermal pressure tracks density with extra shock heating near the current sheet.
    \item Magnetic pressure peaks at flux pile-up in the four current-sheet arms.
  \end{itemize}
\end{frame}

\begin{frame}{Formation of the current sheet}
  \begin{columns}[c]
    \begin{column}{0.33\textwidth}
      \centering
      \includegraphics[width=\linewidth]{rho_t0p100.pdf}\\ \scriptsize t = 0.1
    \end{column}
    \begin{column}{0.33\textwidth}
      \centering
      \includegraphics[width=\linewidth]{rho_t0p300.pdf}\\ \scriptsize t = 0.3
    \end{column}
    \begin{column}{0.33\textwidth}
      \centering
      \includegraphics[width=\linewidth]{rho_final.pdf}\\ \scriptsize t = 0.5
    \end{column}
  \end{columns}
  \vspace{0.3em}
  \begin{itemize}\scriptsize
    \item t = 0.1: still smooth, the vortices are just forming.
    \item t = 0.3: diagonal shocks have formed; bow-tie compression starts.
    \item t = 0.5: fully-developed current sheet and bounding shocks.
  \end{itemize}
\end{frame}

\begin{frame}{Convergence in density}
  \begin{columns}[c]
    \begin{column}{0.55\textwidth}
      \centering
      \includegraphics[width=\linewidth]{convergence.pdf}
    \end{column}
    \begin{column}{0.45\textwidth}
      \centering\scriptsize
      \begin{tabular}{rrr}\toprule
        nx & $L^{1}$ & $L^{2}$ \\\midrule
        32  & 1.26e-1 & 1.73e-1 \\
        64  & 7.27e-2 & 1.07e-1 \\
        128 & 2.87e-2 & 4.26e-2 \\\bottomrule
      \end{tabular}
      \vspace{0.5em}
      \begin{itemize}
        \item $32 \to 64$: rate $\approx 0.79$.
        \item $64 \to 128$: rate $\approx 1.33$.
        \item Between 1st and 2nd order: expected for minmod on a shock-dominated flow~\cite{toro2009}.
      \end{itemize}
    \end{column}
  \end{columns}
\end{frame}

\begin{frame}{Divergence control and conservation}
  \begin{columns}[c]
    \begin{column}{0.55\textwidth}
      \centering
      \includegraphics[width=\linewidth]{divb_history.pdf}
    \end{column}
    \begin{column}{0.45\textwidth}
      \begin{itemize}
        \item GLM~\cite{dedner2002} bounds $\max|\nabla\cdot B|$ once shocks form.
        \item Cell-centered centered differences are a pessimistic diagnostic; the GLM-natural measure is smaller.
        \item Mass, momentum, energy conserved to machine precision ($\sim 10^{-16}$ drift per step) -- see \texttt{run.log}.
      \end{itemize}
    \end{column}
  \end{columns}
\end{frame}

\section{Discussion of Results}

\begin{frame}{Accuracy vs limiter choice}
  \begin{itemize}
    \item Observed density extrema at $128^{2}$: $\rho_{\min}\approx 0.47$, $\rho_{\max}\approx 2.16$.
    \item Athena++~\cite{athena_orszag_tang} with MC limiter at the same resolution: tighter extrema ($\rho_{\min}\sim 0.1$, $\rho_{\max}\sim 3.5$).
    \item Gap is almost entirely the limiter: minmod is the most diffusive TVD limiter and smears current sheets fastest~\cite{toro2009}.
    \item MC or van Leer would close the gap but at a cost in robustness for low-beta cells.
  \end{itemize}
\end{frame}

\begin{frame}{Why between first and second order}
  \begin{itemize}
    \item MUSCL is second-order in smooth regions~\cite{toro2009}.
    \item Any TVD limiter drops to first order at local extrema and at shocks.
    \item Orszag-Tang at t = 0.5 is shock-dominated, so the global rate sits between 1 and 2.
    \item Observed: $\approx 0.8$ at coarse resolutions (shock-dominated), $\approx 1.3$ at fine (smooth regions start to dominate).
  \end{itemize}
\end{frame}

\begin{frame}{Limitations}
  \begin{itemize}
    \item Vectorized NumPy: $128^{2}$ runs in $\sim 7$\,s; $256^{2}$ is the practical ceiling on a laptop.
    \item No constrained-transport option -- GLM is simpler but less strict about $\nabla\cdot B$~\cite{balsara1999}.
    \item 2D Cartesian only; no adaptive refinement.
    \item Reconstruction is linear (MUSCL) -- no WENO, no higher-order time integration.
  \end{itemize}
\end{frame}

\section{Conclusions}

\begin{frame}{Conclusions}
  \begin{itemize}
    \item Built a 2D ideal MHD solver from scratch: Godunov FVM, MUSCL-minmod, HLLD~\cite{miyoshi2005}, GLM~\cite{dedner2002}, SSP-RK2~\cite{shu1988}.
    \item Orszag-Tang at $128^{2}$ reproduces the canonical density/pressure pattern at $t = 0.5$~\cite{londrillo2000,athena_orszag_tang}.
    \item Mass, momentum, energy conserved to machine precision; $\max|\nabla\cdot B|$ bounded by GLM.
    \item Convergence rate between 1 and 2, consistent with MUSCL-minmod on a shock-dominated flow.
  \end{itemize}
  \vspace{0.5em}
  \textbf{Natural extensions.} Constrained transport; resistive MHD; 3D; AMR; Numba/JAX compilation.
\end{frame}

\begin{frame}[allowframebreaks]{References}
  \scriptsize
  \printbibliography[heading=none]
\end{frame}

\end{document}
